%Перечень условных обозначений (при необходимости) оформляется в алфавитном порядке
\sectioncentered*{Перечень условных обозначений}
\addcontentsline{toc}{section}{Перечень условных обозначений}
\label{sec:reduction}

\begin{description}
  \item[БЗ] — база знаний;
  \item[ГИС] — гибридная интеллектуальная система;
  \item[ПрО] — предметная область;
  \item[API] (Application Programming Interface) — программный интерфейс приложения;
  \item[JSON] (JavaScript Object Notation) — текстовый формат обмена данными;
  \item[LLM] (Large Language Model) — большая языковая модель;
  \item[OWL] (Web Ontology Language) — язык описания онтологий;
  \item[RAG] (Retrieval-Augmented Generation) — дополнение генерации извлечением;
  \item[RDF] (Resource Description Framework) — модель представления данных
        в виде троек «субъект — предикат — объект».
\end{description}
